\documentclass[10pt]{article}
\usepackage[margin=20mm]{geometry}
\usepackage{amsmath,amssymb,mathtools,fontspec,unicode-math,xurl,hyperref}
\IfFontExistsTF{DejaVu Serif}{\setmainfont{DejaVu Serif}}{\setmainfont{Latin Modern Roman}}
\setmathfont{latinmodern-math.otf}
\setlength{\parindent}{0pt}\setlength{\parskip}{6pt}\sloppy
\title{Initial scientific assessment of
A universal formula for the \(x-y\) swap in topological recursion}
\author{}
\date{Frozen batch M25-0242; initial review, adjudication pending}
\begin{document}
\maketitle
Authors. Alexander Alexandrov, Boris Bychkov, Petr Dunin-Barkowski, Maxim Kazarian and Sergey Shadrin.\par
Journal identifier. DOI `10.4171/JEMS/1615`; 2025 membership retained from the frozen assignment.\par
Reviewed source. arXiv:2212.00320v3, 30 July 2024, 51 pages. Page references below are both PDF and printed page numbers. The publisher’s version was not obtained for comparison; the findings are restricted to this author manuscript.\par
\section*{Assessment and scope}
Assessment: Minor. Two explicit formula errors occur in the Airy two-point specialization on p. 49. Both admit exact tests within the stated spectral-curve assumptions. They concern a subsidiary application; the inspected dependency chain for the universal swap theorem does not use these formulas. The classification follows this restricted mathematical impact, rather than the ease of a correction. No counterexample to Theorems 1.8 or 1.14 has been established.\par
The review checked the spectral-curve hypotheses, the unstable terms, the finite differential-operator formulation, and the closing pole and period argument. It also independently calculated low-order Airy correlators and tested the suspect two-point identities. The primary executor read pp. 2–4, 6–12, 38–39, 42–43 and 47–49 (the bibliography continuation on p. 49 was not a scientific target). Thirteen source pages were actually viewed as images: 2, 3, 6, 8, 10–12, 39, 42–43 and 47–49. The full Fock-space and graph-combinatorial derivations in Sections 2–4, the entire higher-loop induction, and the other enumerative examples have not been independently certified. Auxiliary checks of specified passages are recorded separately in the provenance.\par
\section*{Hypotheses and positive checks}
The input in Section 1.1 requires simple zeros of \(dx\), regular \(y\) there, and nonvanishing \(dy\) there. In Theorems 1.8 and 1.14, \(x\) is also regular at the simple zeros of \(dy\). Consequently \(dx\) cannot vanish at such a point: that would contradict the original input condition. Thus the disjointness used in Section 5.3 is implied by the stated setting.\par
The initial forms \(\omega_{1,0}^{(0)}=-y\,dx\) and \(\omega_{2,0}^{(0)}=B\) are specified separately. The extra term in (1.13) supplies \(\omega_{0,1}^{(0)}=-x\,dy\), and the two-point initial form remains \(B\) (pp. 8, 42). Stable pole statements must not be applied to this diagonal double pole. The closing argument invokes Proposition 5.3 (p. 39), the cancellation results, and the vanishing of periods via the exact differential in Lemma 5.10 (p. 43). The local exponent \(1+2(s-1)\sum r_i\) in (5.18) is nonnegative, including \(s=1\).\par
For \(S=\mathbb P^1\), \(x=z^2/2\), \(y=z\) and \(B=dz_1dz_2/(z_1-z_2)^2\), the only zero of \(dx\) is the simple zero \(0\); \(dy=dz\) has none. Direct use of (1.1), with involution \(z\mapsto-z\), gives \[\omega_{3,0}^{(0)}=\frac{dz_1dz_2dz_3}{z_1^2z_2^2z_3^2},
 \qquad
 \omega_{1,0}^{(1)}=\frac{dz}{8z^4}.\] These agree with the checked terms of (7.6), (7.10), (7.12) and (7.17). On the swapped side all stable correlators vanish because there are no recursion points. These tests are consistency checks, not a proof of the universal formula.\par
\section*{Finding F1: incorrect differentials in (7.20)}
Location and status. Equation (7.20), p. 49; confirmed in the actual PDF image. With \(x_i=z_i^2/2\), put \(D_i=D_{-x_i}\). On a one-form the definition gives \[D_i(f\,dz_i)=-\partial_{z_i}(f/z_i)\,dz_i,
 \qquad D_i(dx_i)=0.\] For \(k=0\) the left-hand side of (7.20) is \[(D_1+D_2)\frac{dz_1dz_2}{(z_1-z_2)^2}
       =\frac{dz_1dz_2}{z_1^2z_2^2}.\] The printed middle and final expressions, respectively, are \[\frac{dx_1dx_2}{(x_1x_2)^2}
       =\frac{16\,dz_1dz_2}{z_1^3z_2^3},
 \qquad D_1D_2(dx_1dx_2)=0.\] At the admissible distinct points \((z_1,z_2)=(1,2)\) the three coefficients are \(1/4\), \(2\), and \(0\). Thus the displayed equality is false; this is not an ambiguity caused by a pole or by an unstable recursion input.\par
Correction and boundary. The intended identity uses \(z_i\) and \(dz_i\) throughout its right-hand side: \[(D_1+D_2)^{k+1}\frac{dz_1dz_2}{(z_1-z_2)^{2k+2}}
 =\frac{(2k+1)!!\,dz_1dz_2}{(z_1z_2)^{2k+2}}
 =\frac{(D_1D_2)^{k+1}(dz_1dz_2)}{(2k+1)!!}.\] An algebraic derivation is retained in the bounded-check evidence; the primary executor’s symbolic calculations also give zero residual for \(k=0,\ldots,4\). The finite tests alone are not used as a proof for all \(k\). This correction repairs the displayed identity, not every formula in the section automatically.\par
\section*{Finding F2: repeated operator index in (7.19) and (7.21)}
Location and status. The second line of (7.19), and (7.21), p. 49, print \((D_1^3+D_1^3)^i\). The exponential in the preceding line depends on \(w_1^3+w_2^3\), so coefficient extraction instead gives \((D_1^3+D_2^3)^i\).\par
Here is a test at genus one. Write \(F=(z_1-z_2)^{-2}\) and \(G=(z_1-z_2)^{-4}\), suppressing \(dz_1dz_2\) while letting the operators act on its coefficient as above. The corrected expression is \[C=\frac{D_1^3+D_2^3}{24}F+\frac{(D_1+D_2)^2}{4}G
 =\frac{5}{8z_1^6z_2^2}+
   \frac{3}{8z_1^4z_2^4}+\frac{5}{8z_1^2z_2^6}.\] It agrees with a direct residue evaluation of (1.1) using the low-order forms above. At \((1,2)\) it equals \(97/512\). The printed second line, with \(2D_1^3\) in its first term, instead gives \(-13/64\), differing by \(-201/512\). This computation also exposes the loss of symmetry in the two variables.\par
Counter-evidence considered. The preceding exponential and the symmetric polynomial \((u_1^3+u_2^3)^i\) in (7.22) indicate the intended index. They support an index correction, but cannot make the printed line true. Correcting the repeated index in both occurrences removes this particular discrepancy. The \(k=0\) term in the later rewriting with \((D_1+D_2)^{k-1}\) requires its intended cancellation or primitive interpretation; this review does not elevate that presentation issue to a further confirmed defect.\par
\section*{Impact, limitations and disposition}
The dependence proceeds from the universal formula (1.13) to the unramified-\(y\) graph formula (7.9), then to the Airy examples. The two incorrect displayed identities occur at the last stage. They invalidate those literal computational lines and should be corrected before relying on them as written. The initial forms, the low-order residue tests, and the symmetric expression (7.22) provide positive counter-evidence against extending this finding to the main swap theorem.\par
The source manifest’s bounded supplementary-material and code search was retained; no paper-specific code was obtained or executed. Code absence is not a mathematical defect. The arXiv version history and the official publisher page were checked; no equivalence with the journal PDF is asserted. A bounded correction search found no relevant notice, which does not establish that no correction exists.\par
This is an initial scientific assessment. The scientific decision and adjudication remain pending; no human reviews or independent acceptance are claimed. The actual computations, source-page receipts, current compilation inputs, dependencies, logs and report-page checks are bound in the accompanying evidence.\par
9 A. Alexandrov, B. Bychkov, P. Dunin-Barkowski, M. Kazarian and S. Shadrin, A universal formula for the \(x-y\) swap in topological recursion, arXiv:2212.00320v3 (2024), https://arxiv.org/abs/2212.00320v3. EMS Press, article record for DOI 10.4171/JEMS/1615, published 14 April 2025, https://ems.press/journals/jems/articles/14298713.\par
\end{document}
